

Perceptual illusions have long been a topic of fascination to both the general public and scientists alike. The reason for this interest is straightforward: through the course of daily life, visual perception is primarily veridical and objects typically appear as they are.  Yet specific configurations can lead to surprising and counter-intuitive distortion of our perception. An example of such a configuration is the Brentano illusion, depicted in the first row of Figure~\ref{fig:illPlots}. Although the blue and the red lines are the same length, the surrounding arrow heads leads to the illusion than one of the colored lines is longer than the other.

For scientists, visual illusions offer a window into the processes and mechanisms underlying perception \parencite[]{Coren.Girgus.2020}. Recently, individual differences have been focal \parencite[see][]{Tulver.2019}.  Are people who are particularly susceptible to one illusion also more or less susceptible to another illusions?  The hope is that the pattern of correlations that are diagnostic of latent processes or representations.  Specific patterns in correlations may indicate factors or clusters.  For instance, perhaps size-constancy illusions cluster together but are independent of illusions based on the misattribution of 3-D structure.  


There is an active literature on individual differences in susceptibility to various illusions.  Most studies report no common mechanism underlying perceptual illusions \parencite[]{Mazuz.etal.2023, Grzeczkowski.etal.2017, Cretenoud.etal.2019, Bosten.etal.2017, Cottier.etal.2023, Ward.etal.2017}. These studies often find significant correlations within variants of specific illusions, yet they find only weak correlations between different types of illusions (with few exceptions).  The resulting conclusion is that factors influencing susceptibility to illusions are likely multiple and specific rather than singular or general.  These results concord well with a lack of general factors in visual perception studies, where a prevailing explanation is that individual variation likely reflects uncorrelated biological aspects such as retinal structure \parencite[e.g.,][]{Mollon.etal.2017, Bosten.Mollon.2010}.  While this lack of common process characterizes much of the literature, there is a notable  exception.  \textcite{Makowski.etal.2023} finds a single factor that operates across a wide class of illusions.  These authors claim this general factor is positively related to positive personality traits (such aagreeableness and honesty-humility) and negatively related to maladaptive personality traits (such as antagonism, psychoticism, disinhibition, and negative affect). 

There are several methodological difficulties in studying patterns of correlations from individual differences \parencite{Rouder.etal.2023}. We present these difficulties where they are most pronounced, in studies of executive control, and use them to motivate our design and analyses.  The seminal study in executive control is from \textcite{Miyake.etal.2000}, who use individual differences to decompose executive function into processes of inhibition, shifting and updating.  Although this decomposition has been impactful, recent studies have proven far more difficult to interpret.  First, behavioral tasks that purportedly measure the same executive function rarely correlate higher than 0.2 in value \parencite{Rey-Mermet.etal.2018,Rouder.etal.2023,Stahl.etal.2014}. Second, the behavioral tasks have poor test-retest reliability \parencite{Draheim.etal.2019, Hedge.etal.2018} making these low correlation difficult to interpret.  Finally, structural-equation-model analyses of individual difference executive functioning does not replicate even in bootstrapped samples indicating a bias toward publishing underpowered studies\parencite{Karr.etal.2018}.   


These difficulties arise from a statistical issue shown in Figure~\ref{fig:introPlot}.  The left-hand panel shows true scores across a group of participants in two tasks.  The right-hand panel shows observed scores, and these are perturbed by trial noise.  The result of this perturbation is that observed correlations are lower than true values.   This downward bias, called the attenuation of correlation, is well known \parencite{Spearman.1904}.  Even so, it is starkly different than most biases we are accustomed to in experimental psychology.  In much of experimental psychology, we have a large-participant safety net.  If we add more participants, then our estimates become closer to the true value.  This limit does not hold for attenuation of correlation, however.  For the case in Figure~\ref{fig:introPlot}, adding more participants adds more points in the scatter plot, but it does not change the size of the perturbations from trial noise nor does mitigate the attenuation.  In fact, adding more participants results in more confidence in a bad result.   

Researchers occasionally ignore this fact or are, perhaps, unaware of its impact.  They do so by stressing the number of participants in a design without sufficient consideration to trial noise and its problematic effects.  It is routine for researchers to emphasize the number of participants in abstracts and methods without mentioning trial size or trial noise.  It is also routine for researchers make power calculations and publish  confidence intervals that reflect participant variability alone without consideration of attenuation.  All these practices are flawed when there is much trial noise--- they all stress being confident in a wrong answer.  

What to do? \textcite{Matzke.etal.2017} were perhaps the first to propose modeling trial noise along with variability across people and their correlations across tasks in a hierarchical framework.  Rouder and colleagues \parencite[]{Rouder.Haaf.2019, Rouder.etal.2023, Rouder.Mehrvarz.2024} show that correlations are indeed disattenuated in hierarchical models, and \textcite{Mehrvarz.Rouder.2025} show that accurate estimation of these correlations is possible but only under certain conditions and with certain assumptions.  One of the key messages in this work is that some of the usual concepts in individual-difference research, such as reliability, are outdated because they do not take into account trial noise.  And as such, there is a need for new ways of visualizing and communicating results for individual differences.  


We have two goals.  First, we study the nature of individual differences in visual illusions to uncover an underlying factor structure.  Second, we use individual-differences in visual illusions to show how trial noise can and should be incorporated in research.  We provide relatively new measures and ways of plotting data as well as new hierarchical factor models for inference.  


To study individual differences in illusions, we constructed a battery of tasks. In each task, the participant had to adjust an object so that it matched another part of the display.  For example, in the Brentano Illusion, participants adjusted the middle chevron until the size of the blue and red segments appeared equal in length.  We build a pipeline in the analysis section.  To preview, much effort is spent graphically examining how well we could measure each person's bias or score on each task.  It is only after demonstrating that people differ and that we could reliably measure these differences that we perform a latent-variable decomposition of the correlations among the tasks.  To foreshadow, we find a middle position inline with most of the previous research.  First, correlations are small, and this small size is not due to attenuation or statistical issues.  Individual differences in various visual illusions are large, but the correlations across different illusions are truly small.  Second, even though these correlations are small, we know them with relatively high precision, and they form a positive manifold. There is a common factor---susceptibility to illusions.  Unfortunately, it accounts for just a small fraction of the variance.  Much of the variability therefore seems to be idiosyncratic to specific illusions.
